% ch2.tex -- Module 5 ESSENTIALS, Chapter 2: Circular Motion: Forces and Real Situations
% Condensed from the final full book: src/mod5/ch2.tex (same notation, same facts).
% Knowledge points only: no worked examples, no questions, no answers.
% All diagrams are the full book's TikZ code, copied (no figure files are needed).
\chapter[Circular Motion: Forces]{Circular Motion: Forces and Real Situations}\label{ch2:chap}
\covers{M5.IQ2.01, M5.IQ2.02, M5.IQ2.03}

\begin{summarybox}[unbreakable]{In this chapter}
\begin{itemize}
  \item Why moving in a circle at a constant speed is accelerating.
  \item The formulas for $v$, $a_\mathrm{c}$, $F_\mathrm{c}$ and $\omega$.
  \item Which real force provides the centripetal force in each situation.
  \item The whirling-stopper investigation of $F_\mathrm{c}$, $m$, $v$ and $r$.
\end{itemize}
\end{summarybox}

%% ==================================================================
\section{Uniform circular motion and centripetal force}\label{ch2:sec:ucm}
\covers{M5.IQ2.01, M5.IQ2.02, M5.IQ2.03}

\begin{bigidea}
An object moving in a circle at a constant speed is accelerating, because its direction
keeps changing. The acceleration points to the centre. A net force towards the centre,
the centripetal force, causes it.
\end{bigidea}

\begin{keywords}
  \kw{Uniform circular motion}{moving around a circle at a constant speed.}
  \kw{Period ($T$)}{the time for one full turn. Unit: second (s).}
  \kw{Tangent}{a line that touches the circle at one point, at \ang{90} to the radius.}
  \kw{Centripetal acceleration ($a_\mathrm{c}$)}{the acceleration towards the centre of the circle.}
  \kw{Centripetal force ($F_\mathrm{c}$)}{the net force towards the centre of the circle.}
\end{keywords}

\begin{itemize}
  \item Velocity is a \term{vector}: it has a size (the speed) \emph{and} a direction. In a
        circle the direction changes all the time. So the velocity changes: the object
        accelerates.
  \item The velocity is always along the tangent (\cref{ch2:fig:ucm}a). Two velocities a
        moment apart, drawn tail to tail, give a change $\Delta v$ that points to the centre
        (\cref{ch2:fig:ucm}b).
  \item So the acceleration $a_\mathrm{c}$ points to the centre (``centripetal'' means
        ``centre-seeking''). One lap is the \term{circumference} $2\pi r$, done in one
        period $T$, so $v = 2\pi r/T$.
  \item $a_\mathrm{c}$ grows with $v^2$: a faster object turns its velocity arrow more
        quickly, \emph{and} the arrow is longer. At the same speed, a smaller $r$ gives a bigger
        $a_\mathrm{c}$.
  \item By Newton's second law ($\vec{F}_\mathrm{net} = m\vec{a}$), the \term{net force} (the
        total of all the forces) also points to the centre: $F_\mathrm{net} = ma_\mathrm{c} = mv^2/r$.
\end{itemize}

\begin{figure}
  \centering
  \begin{tikzpicture}[x=1cm, y=1cm]
    \begin{scope}
      \draw[hscmuted!55, line width=0.8pt] (0,0) circle (1.4);
      \fill[hscink] (0,0) circle (1.5pt);
      \node[diagramlabel, below, inner sep=3pt] at (0,0) {centre};
      \draw[->, hscmuted, line width=0.6pt] (212:1.75) arc (212:255:1.75);
      \node[diagramlabel, text=hscmuted, anchor=north east] at (232:1.85) {motion};
      \draw[guide] (35:1.4) -- ++(125:1.9);
      \node[diagramlabel, text=hscmuted, anchor=south west, inner sep=1pt] at ($(35:1.4)+(125:1.9)+(-0.1,0.02)$)
        {path if the string breaks};
      \foreach \a in {35, 155, 275} {
        \coordinate (P) at (\a:1.4);
        \draw[velocity] (P) -- ++(\a+90:1.0);
        \draw[accel] (P) -- ($(P)!0.5!(0,0)$);
        \fill[object] (P) circle (4pt);
      }
      \node[diagramlabel, text=hscvel, anchor=west] at ($(35:1.4)+(125:0.9)+(0.08,0)$) {$v$};
      \node[diagramlabel, text=hscacc, anchor=north west, inner sep=1pt] at ($(35:0.85)+(0.05,-0.02)$) {$a_\mathrm{c}$};
      \node[diagramlabel, text=hscvel] at ($(155:1.4)+(245:1.2)$) {$v$};
      \node[diagramlabel, text=hscacc, anchor=south, inner sep=2pt] at ($(155:0.85)+(0.05,0.05)$) {$a_\mathrm{c}$};
      \node[diagramlabel, text=hscvel, anchor=west] at ($(275:1.4)+(5:1.05)$) {$v$};
      \node[diagramlabel, text=hscacc, anchor=west, inner sep=2pt] at ($(275:0.95)+(0.02,-0.1)$) {$a_\mathrm{c}$};
      \node[font=\small\bfseries, text=hscink] at (-1.8,1.8) {(a)};
    \end{scope}
    \begin{scope}[xshift=6.4cm, yshift=-0.8cm]
      \coordinate (O) at (0,0);
      \coordinate (A) at (110:1.8);
      \coordinate (B) at (160:1.8);
      \draw[velocity] (O) -- (A) node[pos=0.55, right, inner sep=2pt, text=hscvel] {$v_1$ (earlier)};
      \draw[velocity] (O) -- (B) node[pos=0.45, below left, inner sep=2pt, text=hscvel] {$v_2$ (later)};
      \draw[accel] (A) -- (B) node[pos=0.5, above left, inner sep=1pt, text=hscacc] {$\Delta v$};
      \fill[hscink] (O) circle (1.3pt);
      \node[font=\small\bfseries, text=hscink] at (-2.1,2.5) {(b)};
    \end{scope}
  \end{tikzpicture}
  \caption{(a) Top view of a ball whirled on a string: $v$ is along the tangent,
    $a_\mathrm{c}$ points to the centre. (b) Two velocities a moment apart, tail to tail:
    $\Delta v$ points to the centre.}
  \label{ch2:fig:ucm}
\end{figure}

\begin{formulacard}{Speed, acceleration and force in a circle}
  \cardformula{v = \frac{2\pi r}{T} \qquad\quad a_\mathrm{c} = \frac{v^2}{r} \qquad\quad
    F_\mathrm{c} = \frac{mv^2}{r}}
  \begin{symbols}
    \sym{$v$}{speed of the object (along the tangent)}{\si{m.s^{-1}}}
    \sym{$r$}{radius of the circle}{\si{m}}
    \sym{$T$}{period: the time for one full turn}{\si{s}}
    \sym{$a_\mathrm{c}$}{centripetal acceleration, towards the centre}{\si{m.s^{-2}}}
    \sym{$F_\mathrm{c}$}{centripetal force: the net force towards the centre}{\si{N}}
    \sym{$m$}{mass of the object}{\si{kg}}
  \end{symbols}
  \whentouse{any object moving in a circle at a constant speed.}
  \trap{$T$ is the period (not tension); use the radius, not the diameter. $F_\mathrm{c}$ is
    the net force: never draw it as an extra arrow.}
\end{formulacard}

\begin{itemize}
  \item $F_\mathrm{c} \propto m$, \ $F_\mathrm{c} \propto v^2$, \ $F_\mathrm{c} \propto 1/r$
        ($\propto$ means ``is proportional to''): double $v$ and $F_\mathrm{c}$ is 4 times
        bigger. With the period, $a_\mathrm{c} = 4\pi^2 r/T^2$ (combine $v = 2\pi r/T$ and
        $a_\mathrm{c} = v^2/r$).
  \item \term{Frequency} $f$ (turns each second, in \si{Hz}): $f = 1/T$ (on the sheet).
        Convert first: \si{km.h^{-1}}~$\div\, 3.6$ gives \si{m.s^{-1}}; revolutions per
        minute (rpm)~$\div\, 60$ gives $f$.
  \item \term{Angular velocity} $\omega$ is the angle turned each second:
        $\omega = \Delta\theta/t$ (on the sheet), in \si{rad.s^{-1}}. One full turn is $2\pi$
        \term{radians} (rad) $= \ang{360}$, so $\omega = 2\pi/T$ (not on the sheet).
  \item The centripetal force is \emph{not} a new kind of force: real forces do this job.
        The real force can be tension (ball on a string), friction (flat bend), the normal
        force (banked track, rotor ride) or gravity (Moon or satellite in orbit).
  \item No force pushes you outwards on a turn. By \term{inertia} (the tendency to keep
        moving in a straight line) your body tries to go straight; the seat, seatbelt and
        door push you \emph{inwards}. ``Centrifugal force'' is not a real force.
  \item If a whirling ball's string breaks, the inward force is gone. By Newton's first law,
        the ball moves in a straight line along the tangent, not straight outwards.
\end{itemize}

%% ==================================================================
\section{Real situations: the forces that provide \texorpdfstring{$F_\mathrm{c}$}{Fc}}\label{ch2:sec:situations}
\covers{M5.IQ2.02, M5.IQ2.03}

\begin{bigidea}
In every situation, find the real forces and split them into a part towards the centre
and a vertical part. The part towards the centre adds up to $mv^2/r$.
\end{bigidea}

\begin{keywords}
  \kw{Tension ($F_\mathrm{T}$)}{the pull of a string, rope or chain, along its length.}
  \kw{Normal force ($F_\mathrm{N}$)}{the push of a surface on an object, at \ang{90} to the surface.}
  \kw{Component}{the part of a force that acts in one chosen direction.}
\end{keywords}

\textbf{Method.} Draw a \term{free-body diagram}: the object with one arrow per real force (no
$F_\mathrm{c}$ arrow).
With $\theta$ from the \emph{vertical}, a force has parts $F\cos\theta$ (vertical) and
$F\sin\theta$ (horizontal). Towards the centre they add up to $mv^2/r$; vertically
(horizontal circle) to zero.

\subsection*{Strings: conical pendulums and swings}
A \term{conical pendulum} is a mass on a string moving in a horizontal circle. Only
$F_\mathrm{T}$ and $mg$ act on it (\cref{ch2:fig:conical}).

\begin{figure}
  \centering
  \begin{tikzpicture}[x=1cm, y=1cm]
    \begin{scope}
      \fill[ground] (-1.5,0) rectangle (2.4,0.22);
      \draw[hscink, line width=0.8pt] (-1.5,0) -- (2.4,0);
      \coordinate (piv) at (0,0);
      \coordinate (B) at (1.548,-2.212);
      \coordinate (C) at (0,-2.212);
      \draw[guide] (piv) -- (0,-2.7);
      \draw[hscmuted, dashed, line width=0.7pt] (C) ellipse (1.548 and 0.36);
      \draw[hscink, line width=1pt] (piv) -- (B) node[pos=0.5, right, diagramlabel] {$L$};
      \draw[hscink, line width=0.5pt] (0,-0.7) arc (270:305:0.7);
      \node[diagramlabel] at (0.19,-0.92) {$\theta$};
      \draw[displacement, <->] (C) -- (B) node[pos=0.5, below, diagramlabel, inner sep=2pt] {$r$};
      \fill[hscink] (C) circle (1.5pt);
      \draw[hscmuted, line width=0.4pt] (C) -- (-1.75,-2.75);
      \node[diagramlabel, anchor=north east, align=right, inner sep=1pt] at (-1.75,-2.75) {centre of\\the circle};
      \fill[object] (B) circle (5pt);
      \fill[hscink] (piv) circle (1.5pt);
      \node[diagramlabel, anchor=east] at (-0.12,-0.25) {pivot};
      \node[font=\small\bfseries, text=hscink] at (-1.3,-0.75) {(a)};
    \end{scope}
    \begin{scope}[xshift=7.4cm, yshift=-1.7cm]
      \coordinate (O) at (0,0);
      \draw[guide] (O) -- (0,1.5);
      \draw[guide] (0,1.5) -- (-1.05,1.5);
      \draw[force] (O) -- (-1.05,1.5) node[pos=0.62, below left, inner sep=1pt, text=hscforce] {$F_\mathrm{T}$};
      \draw[force] (O) -- (0,-1.5) node[pos=0.8, right, text=hscforce] {$mg$};
      \draw[hscink, line width=0.5pt] (0,0.55) arc (90:125:0.55);
      \node[diagramlabel] at (-0.15,0.8) {$\theta$};
      \node[diagramlabel, anchor=west] at (0.05,1.05) {$F_\mathrm{T}\cos\theta$};
      \node[diagramlabel, anchor=south] at (-0.53,1.53) {$F_\mathrm{T}\sin\theta$};
      \draw[displacement] (-1.4,-0.6) -- (-2.5,-0.6);
      \node[diagramlabel, anchor=north, text=hscmuted] at (-1.95,-0.65) {to the centre};
      \fill[object] (O) circle (5pt);
      \node[font=\small\bfseries, text=hscink] at (-2.5,1.8) {(b)};
    \end{scope}
  \end{tikzpicture}
  \caption{A conical pendulum. (a) Side view: the centre of the circle is directly below
    the pivot. (b) Free-body diagram; dashed lines are the components of the tension.}
  \label{ch2:fig:conical}
\end{figure}

\begin{itemize}
  \item The centre of the circle is directly \emph{below} the \term{pivot} (the fixed point),
        so $r = L\sin\theta$. Towards the centre: $F_\mathrm{T}\sin\theta = mv^2/r$.
        Vertically: $F_\mathrm{T}\cos\theta = mg$, so $F_\mathrm{T} = mg/\cos\theta > mg$.
  \item Swing-ride chair (hangs from a spinning frame): $r = d + L\sin\theta$, $d$ = axis to chain top.
  \item A \term{vertical circle} (an upright circle) is \emph{not} uniform: the speed changes.
        At each point the net force to the centre is still $mv^2/r$. Lowest point of a swing:
        $F_\mathrm{T} - mg = mv^2/r$ (you feel heavier). Top, on a string:
        $F_\mathrm{T} + mg = mv^2/r$. Neither is on the sheet.
\end{itemize}

\subsection*{Tilted surfaces: bowls and banked tracks}
\begin{itemize}
  \item A \term{banked track} is higher on the outside of the bend. Its \term{bank angle}
        $\theta$ is its tilt from the horizontal. $F_\mathrm{N}$ (at \ang{90} to the road)
        tilts to the centre, by the same $\theta$ from the vertical.
  \item At the \term{design speed} no friction is needed. Faster: the car tends to slide up
        the slope, so friction acts \emph{down} the slope. Slower: friction acts \emph{up}
        the slope (\cref{ch2:fig:bankfriction}).
\end{itemize}

\begin{figure}
  \centering
  % full book figure ch2:fig:bankfriction (bank angle 20 deg). Added here: the angle theta
  % between F_N and the vertical in (b) (F_N is drawn at 110 deg, so 20 deg from the vertical).
  \begin{tikzpicture}[x=1.2cm, y=1.2cm]
    \foreach \dx in {0, 4.0, 8.0} {
      \begin{scope}[xshift=\dx*1.2cm]
        \fill[ground] (0,0) -- (2.8,1.019) -- (2.8,0.70) -- (0,-0.319) -- cycle;
        \draw[hscink, line width=1pt] (0,0) -- (2.8,1.019);
        \draw[guide] (0,0) -- (0.9,0);
        \draw[hscink, line width=0.5pt] (0.5,0) arc (0:20:0.5);
        \node[diagramlabel, font=\footnotesize] at (0.66,0.1) {$\theta$};
        \begin{scope}[rotate around={20:(1.4,0.510)}]
          \draw[object, rounded corners=1.5pt] (0.98,0.60) rectangle (1.82,0.98);
          \fill[hscink] (1.15,0.60) circle (0.085);
          \fill[hscink] (1.65,0.60) circle (0.085);
        \end{scope}
        \draw[force] (1.314,0.744) -- ++(0,-1.0) node[pos=0.9, right, text=hscforce, inner sep=2pt] {$mg$};
        \draw[force] (1.314,0.744) -- ++(110:1.064) node[pos=1, left, text=hscforce] {$F_\mathrm{N}$};
        \fill[hscink] (1.314,0.744) circle (1.3pt);
      \end{scope}
    }
    % (b): angle of F_N from the vertical
    \draw[guide] (5.314,0.744) -- ++(0,1.1);
    \draw[hscink, line width=0.5pt] (5.314,1.304) arc (90:110:0.56);
    \node[diagramlabel, font=\footnotesize] at ($(5.314,0.744)+(100:0.82)$) {$\theta$};
    % friction: (a) up the slope, (c) down the slope
    \draw[force] (1.314,0.744) -- ++(20:0.75) node[pos=1, above, text=hscforce] {$F_\mathrm{f}$};
    \draw[force] (9.314,0.744) -- ++(200:0.75) node[pos=1, above left, text=hscforce, inner sep=1pt] {$F_\mathrm{f}$};
    \node[diagramlabel, font=\footnotesize, align=center, anchor=north] at (1.4,-0.45)
      {\textbf{(a) slower:} friction up the slope};
    \node[diagramlabel, font=\footnotesize, align=center, anchor=north] at (5.4,-0.45)
      {\textbf{(b) design speed:} no friction};
    \node[diagramlabel, font=\footnotesize, align=center, anchor=north] at (9.4,-0.45)
      {\textbf{(c) faster:} friction down the slope};
  \end{tikzpicture}
  \caption{A car on a banked track, from behind (centre of the bend to the left). Not to
    scale.}\label{ch2:fig:bankfriction}
\end{figure}

\begin{itemize}
  \item At the design speed: $F_\mathrm{N}\sin\theta = mv^2/r$ and $F_\mathrm{N}\cos\theta = mg$,
        so $F_\mathrm{N} = mg/\cos\theta > mg$. Banking means less friction is needed, so wet
        turns are safer.
  \item A marble circling inside a round bowl: the same, with $F_\mathrm{N}$ pointing to the
        bowl's centre. $r$ is the horizontal distance to the centre line, not the bowl's radius.
\end{itemize}

\begin{formulacard}[notonsheet]{Tilted string or surface}
  \cardformula{\tan\theta = \frac{v^2}{rg}}
  \begin{symbols}
    \sym{$\theta$}{angle of $F_\mathrm{T}$ or $F_\mathrm{N}$ from the vertical (bank angle)}{\si{\degree}}
    \sym{$v$}{speed (for a banked road: the design speed)}{\si{m.s^{-1}}}
    \sym{$r$}{radius of the horizontal circle}{\si{m}}
    \sym{$g$}{gravitational acceleration, \qty{9.8}{m.s^{-2}}}{\si{m.s^{-2}}}
  \end{symbols}
  \whentouse{conical pendulum; banked road at its design speed; frictionless bowl.}
  \derive{use the two component equations ($F$ is $F_\mathrm{T}$ or $F_\mathrm{N}$).
    \begin{align*}
      F\sin\theta &= \frac{mv^2}{r} && \text{horizontal: net force towards the centre}\\
      F\cos\theta &= mg && \text{vertical: forces balance}\\
      \tan\theta &= \frac{v^2}{rg} && \text{divide the first line by the second; $F$ and $m$ cancel}
    \end{align*}}
  \trap{for a string, $r$ is not the string length: $r = L\sin\theta$. Do not split forces
    along the slope: the acceleration is horizontal.}
\end{formulacard}

\subsection*{Friction: flat bends and the rotor ride}
\begin{itemize}
  \item \term{Static friction} $F_\mathrm{f}$ stops sliding. It is only as big as needed, up to
        $F_\mathrm{f} \le \mu F_\mathrm{N}$ (not on the sheet). The \term{coefficient of
        friction} $\mu$ (no unit) says how strongly the surfaces grip (about 0.8 on a dry road,
        0.1 on ice).
\end{itemize}

\begin{sidebyside}[0.50]
  \centering
  \begin{tikzpicture}[x=1cm, y=1cm]
    \begin{scope}
      \fill[hscmuted!12] (15:1.7) arc (15:115:1.7) -- (115:2.4) arc (115:15:2.4) -- cycle;
      \draw[hscmuted, line width=0.7pt] (15:1.7) arc (15:115:1.7);
      \draw[hscmuted, line width=0.7pt] (15:2.4) arc (15:115:2.4);
      \fill[hscink] (0,0) circle (1.5pt);
      \node[diagramlabel, anchor=north] at (0,-0.05) {centre};
      \coordinate (Q) at (62:2.05);
      \draw[guide] (0,0) -- (Q);
      \draw[velocity] (Q) -- ++(152:1.0) node[above, text=hscvel] {$v$};
      \draw[force] (Q) -- ($(Q)!0.45!(0,0)$) node[pos=1, below right, inner sep=1pt, text=hscforce] {$F_\mathrm{f}$};
      \draw[object, rotate around={152:(Q)}] ($(Q)+(-0.27,-0.15)$) rectangle ($(Q)+(0.27,0.15)$);
      \node[font=\small\bfseries, text=hscink] at (-1.5,2.3) {(a)};
    \end{scope}
    \begin{scope}[xshift=4.0cm, yshift=1.0cm]
      \coordinate (O) at (0,0);
      \draw[force] (O) -- (0,1.1) node[above, text=hscforce] {$F_\mathrm{N}$};
      \draw[force] (O) -- (0,-1.1) node[below, text=hscforce] {$mg$};
      \draw[force] (O) -- (-1.0,0) node[left, text=hscforce] {$F_\mathrm{f}$};
      \fill[object] (O) circle (5pt);
      \node[font=\small\bfseries, text=hscink] at (-0.9,1.3) {(b)};
    \end{scope}
  \end{tikzpicture}
  \captionof{figure}{A car on a flat bend. (a) Top view. (b) Free-body diagram, seen from
    behind (centre to the left).}\label{ch2:fig:flatbend}
\nextside
  \begin{itemize}
    \item \textbf{Flat bend.} The tyres push the road outwards. By Newton's third law, the
          road pushes the tyres \emph{inwards} (\cref{ch2:fig:flatbend}).
    \item This friction is the centripetal force: $F_\mathrm{f} = mv^2/r$. Vertically,
          $F_\mathrm{N} = mg$.
    \item The road can give at most $\mu mg$. If the bend requires more, the car skids
          towards the outside.
  \end{itemize}
\end{sidebyside}

\begin{sidebyside}[0.31]
  \centering
  \begin{tikzpicture}[x=1cm, y=1cm]
    \draw[hscink, line width=1.2pt] (-1.4,-0.2) -- (-1.4,2.6);
    \draw[hscink, line width=1.2pt] (1.4,-0.2) -- (1.4,2.6);
    \draw[hscmuted, line width=0.7pt] (0,2.6) ellipse (1.4 and 0.25);
    \draw[guide] (0,-0.7) -- (0,3.15);
    \node[diagramlabel, anchor=south] at (0,3.15) {axis};
    \draw[->, hscmuted, line width=0.7pt] (-0.5,3.0) arc (200:340:0.55 and 0.18);
    \fill[ground] (-1.3,-0.75) rectangle (1.3,-0.6);
    \draw[hscink, line width=0.7pt] (-1.3,-0.6) -- (1.3,-0.6);
    \node[diagramlabel, anchor=north, text=hscmuted] at (0,-0.8) {floor has dropped};
    \fill[hscbrand!45] (1.15,0.4) rectangle (1.36,1.6);
    \fill[hscbrand!45] (1.255,1.77) circle (0.15);
    \node[diagramlabel, anchor=west] at (1.45,1.75) {rider};
    \node[diagramlabel, anchor=west] at (1.45,2.4) {wall};
    \draw[displacement, <->] (0,0.6) -- (-1.4,0.6) node[pos=0.5, above, diagramlabel] {$r$};
    \coordinate (R) at (1.255,1.05);
    \draw[force] (R) -- (0.1,1.05) node[pos=0.55, above, text=hscforce] {$F_\mathrm{N}$};
    \draw[force] (R) -- (1.255,2.25) node[pos=0.95, left, text=hscforce] {$F_\mathrm{f}$};
    \draw[force] (R) -- (1.255,-0.15) node[pos=0.9, left, text=hscforce] {$mg$};
    \fill[hscink] (R) circle (1.8pt);
  \end{tikzpicture}
  \captionof{figure}{Forces on a rotor rider.}\label{ch2:fig:rotor}
\nextside
  \begin{itemize}
    \item \textbf{Rotor.} In a \term{rotor} ride a drum spins and the floor drops away. The
          riders stay on the wall (\cref{ch2:fig:rotor}).
    \item The wall's normal force is horizontal, towards the axis. It is the centripetal
          force: $F_\mathrm{N} = mv^2/r$.
    \item Static friction acts up the wall and balances the weight: $F_\mathrm{f} = mg$.
    \item Riders stay up only if $mg \le \mu F_\mathrm{N}$. This sets a slowest safe speed.
          Faster, the friction \emph{limit} grows, but the actual friction stays $mg$.
  \end{itemize}
\end{sidebyside}

\begin{formulacard}[notonsheet, unbreakable]{Friction limits: flat bend and rotor}
  \cardformula{v_\mathrm{max} = \sqrt{\mu g r} \qquad\qquad v_\mathrm{min} = \sqrt{\frac{gr}{\mu}}}
  \begin{symbols}
    \sym{$v_\mathrm{max}$}{fastest speed on a flat bend without skidding}{\si{m.s^{-1}}}
    \sym{$v_\mathrm{min}$}{slowest speed at which rotor riders do not slide down}{\si{m.s^{-1}}}
    \sym{$\mu$}{coefficient of friction (tyres on road; rider on wall)}{no unit}
    \sym{$g$}{gravitational acceleration, \qty{9.8}{m.s^{-2}}}{\si{m.s^{-2}}}
    \sym{$r$}{radius of the bend, or of the drum}{\si{m}}
  \end{symbols}
  \whentouse{a flat bend ($v_\mathrm{max}$); a rotor, floor taken away ($v_\mathrm{min}$).
    Mass makes no difference.}
  \derive{at the limiting speed, friction needed $=$ largest friction. Flat bend:
    $mv_\mathrm{max}^2/r = \mu mg$, so $v_\mathrm{max}^2 = \mu g r$. Rotor:
    $mg = \mu F_\mathrm{N} = \mu m v_\mathrm{min}^2/r$, so $v_\mathrm{min}^2 = gr/\mu$.}
  \trap{$\mu F_\mathrm{N}$ is the \emph{largest} friction, not the actual friction.}
\end{formulacard}

%% ==================================================================
\section{Key experiments}\label{ch2:sec:exp}
\covers{M5.IQ2.01}

\begin{experiment}{Investigating $F_\mathrm{c} = mv^2/r$ with a whirling stopper}
\aim{To investigate how $F_\mathrm{c}$ depends on speed (Part A), radius (Part B) and mass
  (Part C).}
\method{\begin{sidebyside}[0.46]
    \begin{itemize}
      \item Whirl the stopper in a horizontal circle (\cref{ch2:fig:apparatus}). The
            hanging mass stays still, so the tension provides $F_\mathrm{c} = Mg$.
      \item The paper clip keeps the radius $r = L$ fixed.
      \item Time 20 turns (less reaction-time error), three times, and average.
            $T = \text{time} \div 20$; $v = 2\pi L/T$.
    \end{itemize}
  \nextside
    \centering
    \begin{tikzpicture}[x=1cm, y=1cm]
    \draw[hscmuted, dashed, line width=0.7pt] (0,0.5) ellipse (3.0 and 0.32);
    \draw[object] (-0.12,-0.75) rectangle (0.12,0.66);
    \draw[hscink, line width=0.8pt] (0,0.66) -- (3.0,0.5);
    \draw[displacement, <->] (0.05,1.30) -- (3.0,1.14) node[pos=0.5, above, diagramlabel] {$L$};
    \fill[hscink] (3.0,0.5) ellipse (0.16 and 0.13);
    \node[diagramlabel, anchor=south, inner sep=1pt, fill=hscexp!4!white] at (3.0,0.68) {stopper, mass $m$};
    \draw[force] (2.85,0.51) -- (1.95,0.56) node[pos=0.5, below, inner sep=1pt, fill=hscexp!4!white, text=hscforce] {$F_\mathrm{T}$};
    \draw[hscink, line width=0.8pt] (0,-0.75) -- (0,-1.35);
    \draw[hscink, fill=hscmuted!40, line width=0.6pt] (-0.12,-0.98) rectangle (0.12,-0.88);
    \node[diagramlabel, anchor=west] at (0.25,-0.93) {paper clip (keeps $L$ fixed)};
    \node[diagramlabel, anchor=west] at (0.25,-0.42) {tube, held in the hand};
    \draw[object] (-0.33,-1.8) rectangle (0.33,-1.35);
    \node[diagramlabel] at (0,-1.575) {$M$};
    \draw[force] (0,-1.8) -- (0,-2.3) node[right, text=hscforce] {$Mg$};
    \node[diagramlabel, anchor=west] at (0.45,-1.575) {hanging mass};
    \end{tikzpicture}
    \captionof{figure}{Apparatus (side view).}\label{ch2:fig:apparatus}
  \end{sidebyside}
  \begin{itemize}
    \item \textbf{Part A:} change $M$ (so $F_\mathrm{c}$); keep $m$ and $L$ the same. Plot
          $F_\mathrm{c}$ against $v^2$. \textbf{Part B:} change $L$ (keep $M$, $m$); plot $v^2$
          against $r$. \textbf{Part C:} use 1, 2, 3, 4 identical stoppers (keep $M$, $L$); plot
          $v^2$ against $1/m$.
    \item Part A: \term{independent variable} (changed) $M$; \term{dependent variable}
          (measured) the time for 20 turns; \term{controlled variables} (kept the same) $m$,
          $L$, a horizontal circle.
  \end{itemize}}
\results{Part A: $F_\mathrm{c}$ against $v^2$ is a straight line through the \term{origin}
  (0, 0). The gradient of the \term{line of best fit} (the straight line closest to all the
  points) agrees with $m/r$ within 1\,\%. Parts B and C (expected from theory): $v^2$ against
  $r$ and $v^2$ against $1/m$ should also be straight lines through the origin.}
\conclusion{Part A: $F_\mathrm{c} \propto v^2$. This supports $F_\mathrm{c} = mv^2/r$. If the
  Part B and C graphs are straight, then at a fixed speed $F_\mathrm{c} \propto 1/r$ and
  $F_\mathrm{c} \propto m$. ($F_\mathrm{c}$ against $1/T^2$ is also straight; $F_\mathrm{c}$
  against $v$ is a curve.)}
\significance{It confirms $F_\mathrm{c} = mv^2/r$ with a measurable force (a hanging weight).
  Engineers use the same model for bends, banked tracks, rides and satellite orbits.}
\end{experiment}

%% ==================================================================
\begin{summarybox}[unbreakable]{Remember}
\begin{itemize}
  \item Constant speed in a circle is still acceleration, because the direction changes.
        $a_\mathrm{c} = v^2/r$ points to the centre. Also $v = 2\pi r/T$,
        $f = 1/T$ and $\omega = \Delta\theta/t$.
  \item $F_\mathrm{c} = mv^2/r$ is the net force towards the centre, given by a real
        force (tension, normal force, friction, gravity). There is no centrifugal force.
  \item Draw only real forces. Towards the centre they add up to $mv^2/r$; vertically
        (horizontal circle) they add up to zero.
  \item Conical pendulum or banked road (design speed): $\tan\theta = v^2/(rg)$.
        Flat bend: $v_\mathrm{max} = \sqrt{\mu g r}$. Rotor: $v_\mathrm{min} = \sqrt{gr/\mu}$.
  \item Whirling stopper: $F_\mathrm{c} = Mg$; $F_\mathrm{c}$ against $v^2$ is a straight line
        through the origin, with gradient $m/r$.
\end{itemize}
\end{summarybox}
